\documentclass[11pt,a4paper]{article}
\usepackage[utf8]{inputenc}
\usepackage{amsmath,amssymb,amsthm}
\usepackage{listings}
\usepackage{xcolor}
\usepackage[margin=1in]{geometry}
\usepackage{hyperref}

\newtheorem{theorem}{Theorem}

\lstset{
  language=C++,
  basicstyle=\ttfamily\small,
  keywordstyle=\color{blue}\bfseries,
  commentstyle=\color{green!60!black},
  stringstyle=\color{red},
  numbers=left,
  numberstyle=\tiny\color{gray},
  breaklines=true,
  frame=single,
  tabsize=4,
  showstringspaces=false
}

\title{IOI 2005 -- Mountain}
\author{Solution}
\date{}

\begin{document}
\maketitle

\section{Problem Statement Summary}
Maintain a function $f$ on discrete points $\{1, \ldots, N\}$,
initially $f(x) = 0$. Support three operations:
\begin{enumerate}
  \item \textbf{Range add:} given $l, r, v$, set $f(x) \gets f(x) + v$
        for $x \in [l, r]$.
  \item \textbf{Clamp to zero:} set $f(x) \gets \max(f(x), 0)$ for
        all $x$.
  \item \textbf{Threshold query:} given $h$, find the leftmost $x$
        with $f(x) \ge h$, or report that none exists.
\end{enumerate}

\section{Solution: Segment Tree with Lazy Propagation}

\subsection{Lazy Tags}
Each node carries two lazy tags:
\begin{itemize}
  \item \texttt{lazy\_add}: pending additive update.
  \item \texttt{lazy\_set} / \texttt{set\_val}: pending ``set all to
        value'' (from the clamp operation).
\end{itemize}
When pushing down, the ``set'' tag is applied first (overriding
children), then the ``add'' tag.

\subsection{Clamp Operation}
The ``clamp to zero'' is handled recursively:
\begin{itemize}
  \item If $\min \ge 0$ in the segment: no change.
  \item If $\max \le 0$: set the entire segment to 0.
  \item Otherwise: push down lazy tags and recurse into children.
\end{itemize}
This is a form of the \textbf{Segment Tree Beats} technique. Segments
that are entirely non-negative or entirely non-positive are handled in
$O(1)$; only mixed segments require recursion.

\subsection{Threshold Query}
To find the leftmost $x$ with $f(x) \ge h$: if the segment's max
is $< h$, return $-1$. Otherwise, recurse left first; if not found,
recurse right.

\section{C++ Implementation}

\begin{lstlisting}
#include <bits/stdc++.h>
using namespace std;

const int MAXN = 1000005;

struct Node {
    long long mn, mx;
    long long lazy_add;
    bool lazy_set;
    long long set_val;
};

Node tree[4 * MAXN];
int n;

void build(int v, int l, int r) {
    tree[v] = {0, 0, 0, false, 0};
    if (l == r) return;
    int mid = (l + r) / 2;
    build(v * 2, l, mid);
    build(v * 2 + 1, mid + 1, r);
}

void pushDown(int v, int l, int r) {
    if (l == r) return;
    for (int c : {v * 2, v * 2 + 1}) {
        if (tree[v].lazy_set) {
            tree[c].mn = tree[c].mx = tree[v].set_val;
            tree[c].lazy_add = 0;
            tree[c].lazy_set = true;
            tree[c].set_val = tree[v].set_val;
        }
        tree[c].mn += tree[v].lazy_add;
        tree[c].mx += tree[v].lazy_add;
        if (tree[c].lazy_set)
            tree[c].set_val += tree[v].lazy_add;
        else
            tree[c].lazy_add += tree[v].lazy_add;
    }
    tree[v].lazy_add = 0;
    tree[v].lazy_set = false;
}

void pull(int v) {
    tree[v].mn = min(tree[v * 2].mn, tree[v * 2 + 1].mn);
    tree[v].mx = max(tree[v * 2].mx, tree[v * 2 + 1].mx);
}

void rangeAdd(int v, int l, int r, int ql, int qr, long long val) {
    if (ql > r || qr < l) return;
    if (ql <= l && r <= qr) {
        tree[v].mn += val;
        tree[v].mx += val;
        if (tree[v].lazy_set) tree[v].set_val += val;
        else tree[v].lazy_add += val;
        return;
    }
    pushDown(v, l, r);
    int mid = (l + r) / 2;
    rangeAdd(v * 2, l, mid, ql, qr, val);
    rangeAdd(v * 2 + 1, mid + 1, r, ql, qr, val);
    pull(v);
}

void clampZero(int v, int l, int r) {
    if (tree[v].mn >= 0) return;
    if (tree[v].mx <= 0) {
        tree[v].mn = tree[v].mx = 0;
        tree[v].lazy_add = 0;
        tree[v].lazy_set = true;
        tree[v].set_val = 0;
        return;
    }
    if (l == r) {
        tree[v].mn = tree[v].mx = max(tree[v].mn, 0LL);
        return;
    }
    pushDown(v, l, r);
    int mid = (l + r) / 2;
    clampZero(v * 2, l, mid);
    clampZero(v * 2 + 1, mid + 1, r);
    pull(v);
}

int query(int v, int l, int r, int ql, int qr, long long h) {
    if (ql > r || qr < l || tree[v].mx < h) return -1;
    if (l == r) return l;
    pushDown(v, l, r);
    int mid = (l + r) / 2;
    int res = query(v * 2, l, mid, ql, qr, h);
    if (res != -1) return res;
    return query(v * 2 + 1, mid + 1, r, ql, qr, h);
}

int main() {
    ios::sync_with_stdio(false);
    cin.tie(nullptr);

    int Q;
    cin >> n >> Q;
    build(1, 1, n);

    while (Q--) {
        char op;
        cin >> op;
        if (op == 'A') {
            int l, r;
            long long val;
            cin >> l >> r >> val;
            rangeAdd(1, 1, n, l, r, val);
        } else if (op == 'M') {
            clampZero(1, 1, n);
        } else if (op == 'Q') {
            long long h;
            cin >> h;
            cout << query(1, 1, n, 1, n, h) << "\n";
        }
    }

    return 0;
}
\end{lstlisting}

\section{Complexity Analysis}
\begin{itemize}
  \item \textbf{Range add:} $O(\log N)$ per operation.
  \item \textbf{Clamp:} Amortized $O(\log^2 N)$ per operation (Segment
        Tree Beats analysis).
  \item \textbf{Query:} $O(\log N)$ per query.
  \item \textbf{Space:} $O(N)$.
\end{itemize}

\paragraph{Note.} For the full IOI 2005 problem with linear-function
additions ($f(x) \gets f(x) + s + d(x - l)$ on $[l, r]$), each node
stores a lazy tag $(s, d)$ and tracks min/max at segment endpoints. The
clamp logic remains the same.

\end{document}
